\cleardoublepage

\section{绪论}

新视角合成 (Novel View Synthesis) 作为计算机视觉中的一项重要任务，旨在通过图片输入进行场景表示的重建，
从而使得计算机能够从任意视点以任意视角进行图片的渲染。新视角合成在实际应用中也具有巨大的价值，一方面，随着虚拟现实
头戴显示器的快速发展，新视角合成可用于六自由度沉浸式视频的建模，扩展了虚拟现实的能力边界；另一方面，具身智能和智能驾驶
的模型训练需要大量的数据，而在真实场景下进行采集数据具有较大的困难性和危险性。因此，需要一个真实的仿真环境来进行模拟，
而新视角合成的任务则能建立一个高保真的场景模型，从而能够进行高效的仿真模拟。

传统的新视角合成主要是基于点或者面片表示的三维模型，因为它们具有显式的几何约束，并且这样的显式表示能够通过基于GPU/CUDA的渲染器
来实现快速的渲染。然而，基于面片的三维模型由于三角形边的建立不具有可微性，难以从图片得到高质量的三维模型；而基于点云的模型虽然能够
通过如光束平差法\cite{triggs1999bundle}的方法实现稀疏的点云和相机位姿的生成，再通过PMVS\cite{PMVS}等算法得到稠密点云，但由于
点云本身的离散性仍然难以得到较高的渲染质量。

到最近，神经辐射场(neural radiance fields)\cite{mildenhall2021nerf}的方法提出了一种基于多层感知机的连续场景表示，其通过射线步进的体渲染
来进行新视角的合成。朴素的使用多层感知机进行该连续场的建模存在训练较慢、模型参数过多的问题，因此 instantNGP\cite{muller2022instant} 提出用哈希网格进行连续的场景表征，
空间某点的特征向量通过哈希网格的插值得到; plenoxels\cite{fridovich2022plenoxels} 则用体素作为场景表示，空间中某个点的特征向量通过相邻体素特征的插值获得。这些方法
都是将模型构建为三维空间的一个特征向量场，然后再通过多层感知机将特征向量解析成颜色、不透明度等属性，由于方法本身都具有空间连续型的特质，帮助了模型的快速收敛与
优化，同时还达到了更高的渲染质量。

基于神经辐射场的方法通过将三维模型建模为隐式的辐射场，再通过体渲染的方式进行渲染与训练，从而达到了比较优越的渲染质量。然而，体渲染需要借助昂贵的像素采样，这不仅
带来了巨大的渲染负担，同时由于是稀疏采样，也会导致渲染的质量降低。为解决这个问题 kerbl 等人提出了 3d gaussian splatting\cite{kerbl2023} 的算法，将三维场景
通过多个三维高斯分布进行拟合, 并通过基于分块的光栅化器进行快速的渲染，从而在保证较高渲染质量的同时还实现了实时的渲染速度。

3DGS在很多数据集上都达到了SOTA的效果。然而，在实际进行室内场景重建时，3DGS遇到了比较大的困难：要重建出高质量的室内场景时，我们需要对场景的各区域有充分的视角覆盖，
但是用普通的相机要实现该目标会需要比较繁琐的数据采集过程。因此，我们希望能够通过视角更广的相机进行3DGS的重建。然而3DGS的渲染算法核心为 EWA splatting \cite{zwicker2002ewa}，
其对于3DGS到2DGS的投影采用了一个局部的仿射近似，通过一个线性变化对复杂的投影进行拟合。在相机的视场角较窄的情况下，这个近似带来的错误是可以接受的，但在大视场角相机下，
这个近似会带来无法被忽略的误差，因为非线性的畸变会导致真实的投影与对称的二维高斯分布存在显著的偏差，特别是对于在机器人和自动驾驶领域广泛使用的鱼眼相机中。
而本论文提出了一个完整的基于鱼眼相机和雷达的管线，其中包括鱼眼相机数据内外参的标定和怎么通过鱼眼图片和雷达点云进行高质量3DGS的训练。

本论文的目标是能够从多张鱼眼图片和雷达数据中重建出一个高质量的三维高斯模型，从而能通过该模型得到进行高质量的新视角合成结果, 并且在此基础上实现一个效率较高的优化方案。在鱼眼数据的三维高斯重建上，
最近的方法如3DGUT\cite{wu20253dgut}和3DGEER\cite{huang20253dgeer}在这个方面进行了尝试，3DGUT通过抽样近似的方式，把三维高斯在鱼眼上的投影近似为一个二维的椭圆，椭圆的形状为特定的若干
控制点拟合而成，这样的做法实现了近似的鱼眼投影拟合，但是对于我们使用的视场角超过180度的相机，其会在边缘部分会带来较为严重的错误；3DGEER则是通过射线采样的方式进行训练，通过逐射线的渲染使得渲染过程可以
适配任意的畸变模型，然而这样逐射线的训练的效率过低。

我们对该部分的解法主要基于三个关键模块。我们首先引入雷达辅助的相机运动结构恢复(Structure from Motion)管线来加速我们相机内外参的优化过程。由于场景注册需要的图片数量过多，传统colmap中逐图片注册并在特定间隔后进行
全局光束平差优化的运动结构恢复管线的效率过低，特别是将相机内外参进行共同优化时，全局光束平差的时间消耗过长，在SfM注册图片数量过多时尤为如此。因此，我们通过FAST-LIO2\cite{xu2022fastlio}的方法，该方法以惯性测量单元（IMU）
和激光雷达（LiDAR）得到的原始点云作为输入，同时估计传感器的运动轨迹和融合后的点云。然后，我们再通过边缘对齐的方法\cite{yuan2021pixel}进行激光雷达与相机之间外参的标定。于是，我们就可以在完成数据采集后就得到一个较精确到粗外参。
我们还通过ChArUco标定版进行了鱼眼相机内参的预标定，这样的方法为我们的SfM提供了一个良好的初始值，从而解决了传统全局光束平差法容易遇到的局部最小值的问题。第二个模块则是我们的透视图采样法。直接通过原始的鱼眼图片进行三维高斯的训练会面临
光栅化形状不规则，精确拟合计算成本较高的问题，而通过透视图进行三维高斯的训练则已经被验证能得到较良好的效果，并且存在较多可行的优化方案。在这样的前提下，我们认为将鱼眼相机重建回归到透视相机的重建会是比较可行的方案，于是我们设计了一个
鱼眼图片拆分算法，可以将单张鱼眼图片拆分为多张等价的透视图进行训练，在保证训练时能利用鱼眼高视场角的丰富信息的同时，能够以较高的效率进行三维高斯的训练。第三个模块则是基于激光雷达的深度约束优化，传统三维高斯在弱纹理的区域会出现大量floater，
导致在新视角合成时出现明显的伪影问题。我们的算法则以稀疏点云和鱼眼中生成的透视图作为输入，通过高质量的预训练深度预测大模型进行稠密深度的预测，然后通过截断距离场构建出深度约束网格，在进行三维高斯优化时进行对应的法线和深度约束，从而提高新视角合成的质量。

总而言之，本文的主要贡献可概括为以下三点：

\begin{enumerate}
    \item 提出了一种激光雷达与惯性测量单元辅助的快速运动结构恢复管线，用于提升鱼眼相机数据内外参优化与场景注册的效率和稳定性。
    \item 提出了一种基于透视图采样的鱼眼图像三维高斯重建方法，将鱼眼图像拆分为多张等价透视图，在保留大视场角信息的同时提高了模型训练效率。
    \item 提出了一种从稀疏激光雷达点云到深度约束网格的构建方法，并将法线与深度约束引入三维高斯优化过程，从而提升弱纹理区域的新视角合成质量。
\end{enumerate}

我们在公开数据集和自己采集的数据上的结果都展现了我们能通过鱼眼相机和激光雷达重建得到比之前的方法更高质量三维模型。同时，在保证质量的同时，我们的方法也实现了比之前的算法更快的训练速度。
